\documentclass[11pt,a4paper]{article}

% Codificacao e idioma
\usepackage[T1]{fontenc}
\usepackage[utf8]{inputenc}
\usepackage[brazil]{babel}
\usepackage{lmodern}
\usepackage{microtype}

% Pagina e elementos graficos
\usepackage[a4paper,margin=2.5cm,headheight=15pt]{geometry}
\usepackage{xcolor}
\usepackage{enumitem}
\usepackage{titlesec}
\usepackage{fancyhdr}
\usepackage{hyperref}
\usepackage{bookmark}
\usepackage{tcolorbox}
\usepackage{lastpage}

\definecolor{navy}{HTML}{17324D}
\definecolor{teal}{HTML}{168C8C}
\definecolor{ink}{HTML}{263442}
\definecolor{muted}{HTML}{657586}
\definecolor{pale}{HTML}{EAF4F4}
\definecolor{palegray}{HTML}{F3F6F8}

\hypersetup{
  colorlinks=true,
  linkcolor=navy,
  urlcolor=teal,
  pdftitle={Sistema de Gestao de Drogaria — Fase 1},
  pdfsubject={Levantamento e especificacao de requisitos},
  pdfauthor={Projeto de Banco de Dados}
}

\setlength{\parindent}{0pt}
\setlength{\parskip}{5pt}
\setlist[itemize]{leftmargin=1.35em,itemsep=2pt,topsep=3pt}
\setlist[enumerate]{leftmargin=1.7em,itemsep=3pt,topsep=4pt}

\titleformat{\section}{\Large\bfseries\color{navy}}{\thesection.}{0.55em}{}
\titleformat{\subsection}{\large\bfseries\color{teal}}{}{0pt}{}
\titlespacing*{\section}{0pt}{20pt}{7pt}
\titlespacing*{\subsection}{0pt}{12pt}{4pt}

\pagestyle{fancy}
\fancyhf{}
\fancyhead[L]{\small\color{muted}Projeto de Banco de Dados}
\fancyhead[R]{\small\color{muted}Sistema de Gestao de Drogaria}
\fancyfoot[L]{\small\color{muted}Fase 1 — Requisitos}
\fancyfoot[R]{\small\color{muted}\thepage\ de \pageref{LastPage}}
\renewcommand{\headrulewidth}{0.4pt}
\renewcommand{\footrulewidth}{0pt}

\newcommand{\req}[2]{\item[\textbf{#1}] #2}
\newcommand{\ruleitem}[2]{\item[\textbf{#1}] #2}

\begin{document}

\begin{titlepage}
\thispagestyle{empty}
\vspace*{2.2cm}
{\color{teal}\large PROJETO DE BANCO DE DADOS\par}
\vspace{1.2cm}
{\Huge\bfseries\color{navy}Sistema de Gestao de Drogaria\par}
\vspace{0.45cm}
{\Large\color{ink}Fase 1: Levantamento e especificacao de requisitos\par}
\vspace{1.2cm}
\begin{tcolorbox}[colback=pale,colframe=pale,boxrule=0pt,left=12pt,right=12pt,top=10pt,bottom=10pt]
\textbf{Disciplina}\\
SPOBDD2 — Elaborar em fases o Plano do Projeto de BD
\end{tcolorbox}
\vfill
{\large\color{ink}Outubro de 2026\par}
\vspace{1.5cm}
{\small\color{muted}Documento de requisitos do projeto\par}
\end{titlepage}

\tableofcontents
\clearpage

\section{Identificacao do projeto}
\textbf{Projeto:} Sistema de Gestao de Drogaria\\
\textbf{Disciplina:} SPOBDD2 — Elaborar em fases o Plano do Projeto de BD

\section{Descricao do sistema}
O projeto propoe um sistema de banco de dados para apoiar as operacoes de uma drogaria, abrangendo o cadastro de produtos e clientes, controle de estoque por lote quando aplicavel, fornecedores e compras, vendas, pagamentos, receitas medicas digitalizadas ou digitais e controle de acesso por perfil de usuario.

O catalogo podera conter medicamentos e produtos nao medicamentosos, como itens de higiene pessoal, fraldas, cosmeticos e outros produtos comercializados pela drogaria. O sistema devera manter o historico relevante das operacoes, possibilitando consultas e rastreabilidade.

\section{Objetivo}
Organizar e manter consistentes os dados operacionais da drogaria, permitindo registrar compras e recebimentos, acompanhar estoque, registrar vendas e pagamentos, controlar cancelamentos e associar receitas as vendas quando aplicavel.

\section{Escopo do projeto}
\subsection{Incluido no escopo}
\begin{itemize}
\item Cadastro e manutencao de produtos, categorias e fabricantes.
\item Controle de estoque e de lotes, incluindo datas de validade quando aplicaveis.
\item Cadastro de fornecedores e multiplos contatos por fornecedor.
\item Pedidos de compra separados do recebimento efetivo das mercadorias.
\item Cadastro de clientes e funcionarios.
\item Vendas com itens, descontos percentuais por item e associacao aos lotes dispensados.
\item Pagamentos em multiplas modalidades, pagamentos parciais, saldo pendente e registro de estornos.
\item Cadastro de receitas digitais ou digitalizadas por referencia a documento externo, com controle de acesso e associacao a uma unica venda no escopo deste projeto.
\item Perfis de acesso iniciais: atendente, gerente e administrador.
\item Cancelamento integral de vendas, com tratamento separado da reversao de estoque e dos estornos financeiros.
\item Controle de estoque por movimentacoes rastreaveis, com identificacao dos lotes utilizados nas vendas e recebidos nas compras.
\item Cadastro de medicamentos com multiplos principios ativos, por meio de cadastro proprio de principios ativos e relacionamento associativo.
\item Consultas operacionais para apoiar o acompanhamento de vendas, estoque, compras e pagamentos.
\end{itemize}

\subsection{Fora do escopo nesta fase}
\begin{itemize}
\item Construcao de interface completa ou aplicacao pronta para producao.
\item Implementacao detalhada de todas as regras sanitarias e regulatorias para cada classe de medicamento.
\item Controle de dispensacoes de uma mesma receita em multiplas vendas.
\item Armazenamento dos arquivos de receita diretamente no banco de dados; sera mantida uma URL ou referencia protegida para documento armazenado externamente.
\item Gestao avancada de convenios, planos de saude, entregas, programa de fidelidade ou integracao com sistemas externos.
\item Criacao de uma entidade/tabela de relatorios; os relatorios previstos serao consultas SQL e poderao, se conveniente, ser implementados como views.
\item Devolucoes parciais de produtos; nesta fase, sera contemplado apenas o cancelamento integral da venda.
\end{itemize}

\section{Requisitos funcionais}
\subsection{Produtos e estoque}
\begin{description}[style=nextline,leftmargin=1.2cm,labelsep=0.25cm]
\req{RF01}{O sistema deve permitir cadastrar, consultar, atualizar e inativar produtos comercializados pela drogaria, incluindo medicamentos e produtos nao medicamentosos.}
\req{RF02}{O sistema deve permitir classificar produtos por categoria e registrar informacoes de fabricante quando aplicavel.}
\req{RF03}{O sistema deve controlar o estoque dos produtos e, quando aplicavel, registrar lotes, quantidades e datas de validade.}
\req{RF04}{O sistema deve permitir consultar a quantidade disponivel por produto e, quando houver controle por lote, por lote.}
\req{RF05}{O sistema deve identificar lotes vencidos e lotes proximos do vencimento, para apoiar o controle do estoque.}
\req{RF06}{O sistema deve registrar movimentacoes de estoque decorrentes de recebimentos, vendas, cancelamentos e outros ajustes autorizados, mantendo historico da operacao.}
\end{description}

\subsection{Fornecedores e compras}
\begin{description}[style=nextline,leftmargin=1.2cm,labelsep=0.25cm]
\req{RF07}{O sistema deve permitir cadastrar, consultar, atualizar e inativar fornecedores, armazenando seus dados cadastrais basicos, incluindo razao social, CNPJ e endereco.}
\req{RF08}{O sistema deve permitir registrar varios telefones e e-mails para um mesmo fornecedor, mantendo cada contato associado ao fornecedor correspondente.}
\req{RF09}{O sistema deve permitir registrar pedidos de compra associados a um fornecedor, incluindo data, itens solicitados, quantidades e precos unitarios negociados.}
\req{RF10}{O sistema deve permitir registrar o recebimento de mercadorias separadamente do pedido de compra, incluindo data, quantidades efetivamente recebidas e, quando aplicavel, lote e validade.}
\req{RF11}{O sistema deve permitir recebimentos parciais de um pedido e acompanhar se o pedido esta pendente, parcialmente recebido ou concluido.}
\end{description}

\subsection{Clientes, funcionarios e acesso}
\begin{description}[style=nextline,leftmargin=1.2cm,labelsep=0.25cm]
\req{RF12}{O sistema deve permitir cadastrar, consultar, atualizar e inativar clientes, preservando o historico de vendas.}
\req{RF13}{O sistema deve permitir cadastrar funcionarios e associar cada usuario a um perfil de acesso.}
\req{RF14}{O sistema deve oferecer perfis iniciais de atendente, gerente e administrador, com permissoes distintas para as funcionalidades.}
\req{RF15}{O sistema deve permitir que usuarios autorizados consultem dados cadastrais e operacionais conforme suas permissoes, restringindo especialmente o acesso a receitas e dados associados.}
\end{description}

\subsection{Vendas e cancelamentos}
\begin{description}[style=nextline,leftmargin=1.2cm,labelsep=0.25cm]
\req{RF16}{O sistema deve permitir registrar vendas com data e hora, funcionario responsavel, cliente quando aplicavel, itens, quantidades e valores.}
\req{RF17}{O sistema deve permitir aplicar desconto percentual individualmente a cada item da venda e calcular os valores dos itens e o total da venda.}
\req{RF18}{O sistema deve verificar a disponibilidade de estoque antes de confirmar uma venda e impedir a dispensacao de lotes vencidos.}
\req{RF19}{Quando houver controle por lote, o sistema deve registrar quais lotes foram utilizados para atender cada item vendido.}
\req{RF20}{O sistema deve permitir cancelar vendas somente a gerente ou a funcionario com permissao especifica, exigindo motivo e registrando data, hora e responsavel.}
\req{RF21}{Ao cancelar uma venda, o sistema deve reverter de forma consistente os efeitos no estoque e preservar o registro original da venda.}
\end{description}

\subsection{Pagamentos e estornos}
\begin{description}[style=nextline,leftmargin=1.2cm,labelsep=0.25cm]
\req{RF22}{O sistema deve permitir registrar um ou mais pagamentos associados a mesma venda, com valor, data e forma de pagamento.}
\req{RF23}{O sistema deve contemplar Pix, dinheiro, cartao de debito e cartao de credito, permitindo combinar modalidades na mesma venda.}
\req{RF24}{O sistema deve permitir pagamentos parciais, calcular o saldo pendente e indicar se a venda esta pendente, parcialmente paga ou quitada.}
\req{RF25}{O sistema deve permitir registrar estornos associados aos pagamentos originais, preservando o historico e registrando valor, data e motivo.}
\end{description}

\subsection{Receitas}
\begin{description}[style=nextline,leftmargin=1.2cm,labelsep=0.25cm]
\req{RF26}{O sistema deve permitir registrar dados essenciais de receitas digitais ou digitalizadas, incluindo datas pertinentes, nome e registro profissional do prescritor e referencia/URL para o documento.}
\req{RF27}{O sistema deve permitir associar cada receita a uma unica venda no escopo deste projeto.}
\req{RF28}{O sistema deve registrar, quando aplicavel, os medicamentos e as quantidades prescritas e dispensadas associados a receita e a venda correspondente.}
\req{RF29}{O sistema deve restringir o acesso as receitas e aos documentos referenciados a usuarios autorizados.}
\end{description}

\subsection{Consultas e relatorios operacionais}
\begin{description}[style=nextline,leftmargin=1.2cm,labelsep=0.25cm]
\req{RF30}{O sistema deve permitir consultar vendas por periodo e funcionario, incluindo quantidade de vendas e total vendido, desconsiderando vendas canceladas nos totais operacionais.}
\req{RF31}{O sistema deve permitir consultar produtos com estoque abaixo de um limite minimo definido para o produto, caso esse limite seja cadastrado.}
\req{RF32}{O sistema deve permitir consultar lotes vencidos e proximos do vencimento.}
\req{RF33}{O sistema deve permitir consultar compras por fornecedor.}
\req{RF34}{O sistema deve permitir consultar vendas com pagamentos pendentes ou saldo em aberto.}
\end{description}

\section{Requisitos nao funcionais}
\begin{description}[style=nextline,leftmargin=1.2cm,labelsep=0.25cm]
\req{RNF01}{\textbf{Integridade dos dados:} o banco deve preservar relacionamentos consistentes e impedir referencias invalidas entre registros.}
\req{RNF02}{\textbf{Seguranca:} o acesso as funcionalidades deve respeitar a autenticacao e as permissoes atribuidas a cada perfil.}
\req{RNF03}{\textbf{Privacidade:} dados pessoais e informacoes de receitas devem ter acesso restrito e ser tratados conforme os principios e obrigacoes aplicaveis da LGPD.}
\req{RNF04}{\textbf{Consistencia transacional:} operacoes relacionadas, como confirmar uma venda e atualizar o estoque, devem ser executadas de forma consistente, evitando que apenas parte da operacao seja efetivada.}
\req{RNF05}{\textbf{Rastreabilidade:} o sistema deve manter historico das movimentacoes de estoque, vendas, cancelamentos, pagamentos e estornos relevantes.}
\req{RNF06}{\textbf{Backup e recuperacao:} deve existir procedimento para copia de seguranca e recuperacao dos dados em caso de falha.}
\req{RNF07}{\textbf{Manutenibilidade:} a organizacao dos dados deve permitir evolucao futura do sistema sem exigir reconstrucao integral do banco.}
\end{description}

\section{Regras de negocio}
\begin{description}[style=nextline,leftmargin=1.2cm,labelsep=0.25cm]
\ruleitem{RB01}{Todo pedido de compra deve estar associado a um fornecedor e conter pelo menos um item.}
\ruleitem{RB02}{As quantidades solicitadas e os precos unitarios de itens de compra devem ser maiores que zero.}
\ruleitem{RB03}{Registrar um pedido de compra nao aumenta o estoque; somente as quantidades efetivamente recebidas podem aumentar o estoque.}
\ruleitem{RB04}{O recebimento de um item nao pode exceder a quantidade solicitada no pedido correspondente. Se for necessaria quantidade adicional, o pedido devera ser ajustado previamente ou devera ser criado outro pedido.}
\ruleitem{RB05}{Um pedido pode ser recebido parcialmente e deve ter situacao compativel com o que ja foi recebido.}
\ruleitem{RB06}{Um fornecedor pode possuir varios telefones e varios e-mails, registrados separadamente e associados a ele.}
\ruleitem{RB07}{Uma venda deve conter pelo menos um item e as quantidades dos itens devem ser maiores que zero.}
\ruleitem{RB08}{A identificacao do cliente e opcional em vendas comuns, mas obrigatoria quando exigida pelas regras aplicaveis ou pela operacao.}
\ruleitem{RB09}{Cada venda deve identificar o funcionario responsavel pelo registro.}
\ruleitem{RB10}{A quantidade vendida nao pode exceder a quantidade disponivel no estoque elegivel para venda.}
\ruleitem{RB11}{Lotes vencidos nao podem ser utilizados para dispensacao.}
\ruleitem{RB12}{Quando houver controle por lote, cada lote deve estar associado ao produto correspondente e ser rastreavel nas movimentacoes pertinentes.}
\ruleitem{RB13}{O preco unitario e o desconto aplicado devem ser preservados no item da venda, mesmo que o preco de cadastro do produto seja alterado posteriormente.}
\ruleitem{RB14}{O desconto percentual de cada item deve estar entre 0\% e 100\% e nao pode resultar em valor final negativo.}
\ruleitem{RB15}{Somente gerente ou funcionario com permissao especifica pode cancelar uma venda; o motivo do cancelamento e obrigatorio.}
\ruleitem{RB16}{O cancelamento nao deve apagar a venda original e deve reverter os efeitos de estoque de forma consistente.}
\ruleitem{RB17}{Uma venda pode ter varios pagamentos e modalidades, incluindo Pix, dinheiro, cartao de debito e cartao de credito.}
\ruleitem{RB18}{O saldo pendente corresponde ao total devido menos os pagamentos validos, considerados os estornos registrados.}
\ruleitem{RB19}{Os pagamentos validos nao podem exceder o saldo devido, exceto se uma regra explicita de troco ou ajuste financeiro for posteriormente incluida no escopo.}
\ruleitem{RB20}{Um estorno deve referenciar o pagamento original, preservar o historico e nao exceder o valor ainda nao estornado desse pagamento.}
\ruleitem{RB21}{O cancelamento de uma venda e o estorno financeiro sao operacoes distintas; cancelar uma venda nao deve apagar os pagamentos ja registrados.}
\ruleitem{RB22}{No escopo deste projeto, cada receita sera associada a uma unica venda; o sistema nao controlara a utilizacao da mesma receita em vendas sucessivas.}
\ruleitem{RB23}{A referencia ao documento da receita deve apontar para um recurso protegido, e nao para um endereco publico irrestrito.}
\ruleitem{RB24}{A validade e os requisitos aplicaveis a receita e ao medicamento devem ser respeitados; uma imagem digitalizada, por si so, nao garante validade legal.}
\ruleitem{RB25}{Produtos inativados nao devem ser usados em novas vendas, mas seus registros historicos devem permanecer disponiveis para consulta autorizada.}
\ruleitem{RB26}{O acesso as funcionalidades deve obedecer ao perfil e as permissoes do usuario; dados de receitas devem ter protecao adicional conforme a necessidade de acesso.}
\end{description}

\subsection{Regras complementares consolidadas para orientar o DER conceitual}
\begin{description}[style=nextline,leftmargin=1.2cm,labelsep=0.25cm]
\ruleitem{RB27}{Cada produto deve estar associado a exatamente uma categoria principal, sem hierarquia de subcategorias, e a um fabricante. O fabricante nao se confunde com o fornecedor.}
\ruleitem{RB28}{Um medicamento e um tipo especializado de produto. Um medicamento pode possuir um ou varios principios ativos; a associacao sera representada por uma entidade associativa entre Medicamento e Principio Ativo.}
\ruleitem{RB29}{O codigo interno do produto e obrigatorio e unico. O codigo de barras e opcional e, quando informado, deve ser unico.}
\ruleitem{RB30}{O preco atual de venda fica cadastrado em Produto. O preco unitario e o desconto praticados devem ser preservados em cada Item de Venda. O custo negociado de aquisicao fica em Item de Compra.}
\ruleitem{RB31}{O fornecedor e associado a Compra, e nao e fixado diretamente no cadastro do produto. Uma compra pertence a um fornecedor e pode conter varios itens.}
\ruleitem{RB32}{A identificacao do lote e unica pela combinacao Produto + numero do lote. A data de validade pode ficar ausente quando nao for aplicavel ao produto.}
\ruleitem{RB33}{Quando houver datas de validade, a separacao de produtos para venda deve priorizar os lotes elegiveis com vencimento mais proximo (FEFO). Lotes vencidos ou nao elegiveis nao podem ser dispensados.}
\ruleitem{RB34}{O saldo de estoque sera apurado a partir das movimentacoes registradas, sem manter uma quantidade independente como fonte oficial no cadastro do lote. O estoque disponivel nao pode ficar negativo.}
\ruleitem{RB35}{Cada Item de Recebimento deve gerar uma movimentacao de entrada vinculada ao lote correspondente. As movimentacoes de saida de venda devem permitir identificar o Item de Venda e os lotes consumidos.}
\ruleitem{RB36}{Cada Ajuste de Estoque deve registrar motivo, data e funcionario responsavel. Um ajuste pode afetar varios lotes, gerando uma movimentacao para cada lote afetado.}
\ruleitem{RB37}{O cancelamento integral deve preservar a venda original e identificar o funcionario responsavel. Os produtos so devem retornar ao estoque vendavel nos lotes originais quando estiverem aptos; itens improprios nao devem aumentar o estoque vendavel.}
\ruleitem{RB38}{A situacao operacional da venda sera Em andamento, Concluida ou Cancelada. A venda somente podera ser marcada como Concluida apos a confirmacao do pagamento integral; a situacao financeira sera apurada com base nos itens, pagamentos validos e estornos.}
\ruleitem{RB39}{Cada venda, ajuste de estoque e cancelamento deve identificar um funcionario responsavel. Funcionarios e fabricantes/fornecedores podem ser inativados, preservando registros historicos relacionados.}
\ruleitem{RB40}{Telefone, e-mail e endereco do cliente serao modelados em entidades separadas. Telefone e e-mail do fornecedor tambem serao entidades separadas, permitindo multiplos contatos por cadastro.}
\ruleitem{RB41}{Cada funcionario pode ter zero ou uma conta de usuario; cada usuario pertence a exatamente um funcionario e possui um perfil de acesso principal.}
\ruleitem{RB42}{Uma venda pode estar associada a zero ou varias receitas; cada receita pertence a exatamente uma venda e possui um prescritor. Os itens de receita so podem referenciar produtos cadastrados como medicamentos.}
\ruleitem{RB43}{Um ajuste de estoque pode gerar varias movimentacoes, uma por lote afetado. Cada movimentacao deve manter uma origem compativel com sua finalidade operacional, permitindo rastreabilidade entre recebimento, venda, cancelamento ou ajuste.}
\end{description}

\section{Premissas e observacoes}
\begin{itemize}
\item O banco de dados sera desenvolvido em MySQL, com modelagem e implementacao previstas para fases posteriores.
\item A URL ou referencia da receita aponta para documento armazenado fora do banco; o armazenamento externo devera oferecer controle de acesso apropriado.
\item As regras de validade, assinatura e dispensacao de receitas variam conforme o tipo de documento e medicamento. O sistema devera respeitar os requisitos aplicaveis, sem pretender substituir uma analise regulatoria completa.
\item Os relatorios operacionais serao implementados como consultas SQL e, se util, views. Nao sera criada uma entidade de negocio chamada Relatorio.
\item Esta fase define requisitos; DER, modelo logico, normalizacao, script SQL, criacao de tabelas e testes serao desenvolvidos nas fases seguintes.
\end{itemize}

\section{Conclusao}
O levantamento estabelece o escopo funcional inicial do sistema de gestao de drogaria e registra as regras de integridade, rastreabilidade, seguranca e operacao consolidadas. Esta especificacao servira de base para o DER conceitual em notacao de Peter Chen, seguido do modelo logico, da analise de normalizacao e da implementacao fisica no MySQL.

\end{document}
